\writeSectionFrame{\presentationTitle}{Chain of Responsibility}
\begin{frame}{Steckbrief}
	\begin{itemize}
 		\item Deutscher Name: Zuständigkeitskette
 		\item Englischer Name: Chain of Responsibility
 		\item Gruppe: Verhaltensmuster
 	\end{itemize}
\end{frame}

\begin{frame}{Zuständigkeitskette}
	\begin{block}{Zweck}
 		Vermeide die Kopplung des Auslösers einer Anfrage an seinen Empfänger, indem mehr als ein Objekt die Möglichkeit erhält, die Anfrage zu erledigen. Verkette die empfangenden Objekte und leite die Anfrage an der Kette entlang, bis ein Objekt sie erledigt.
 	\end{block}
\end{frame}

\begin{frame}{Allgemeines Klassendiagramm}
    \begin{tikzpicture}[scale=0.7, transform shape]
        % Define the classes
        \umlclass[x=-3,y=0]{Client}{
        }{ 
        }
        
        \umlclass[x=4,y=0]{Handler}{
          + next : Handler \\
        }{ 
          + handleRequest() : void \\
        }
        
        \umlclass[x=1,y=-4]{ConcreteHandler1}{
        }{ 
          + handleRequest() : void \\
        }
        
        \umlclass[x=7,y=-4]{ConcreteHandler2}{
        }{ 
          + handleRequest() : void \\
        }
        
        \umlinherit{ConcreteHandler1}{Handler}
        \umlinherit{ConcreteHandler2}{Handler}
        
        % Relationships between handlers
        \umluniassoc{Client}{Handler}
        \umluniassoc[arg=Nachfolger, angle1=-90, angle2=-149, loopsize=2cm]{Handler}{Handler}
    \end{tikzpicture}
\end{frame}

\frameWithImageOnly{\BILDERPFAD/chainOfResponsibility1.jpg}
\note{
    Ein Kundendienstsystem hat verschiedene Support-Level. Wenn eine Anfrage eingeht, wird sie zuerst an den Frontline-Support gesendet. Falls der Frontline-Support die Anfrage nicht bearbeiten kann, wird sie an den höheren Support-Level (z.B. technischen Support oder Manager) weitergeleitet.
}

\begin{frame}{Klassendiagramm Support}
    \begin{tikzpicture}[scale=0.7, transform shape]
        % Define the classes
        \umlclass[x=-3,y=0]{Client}{
        }{ 
        }
        
        \umlclass[x=4,y=0]{SupportHandler}{
        }{ 
          + handleRequest(issueType: String) : void \\
        }
        
        \umlclass[x=1,y=-4]{BillingSupportHandler}{
        }{ 
          + handleRequest(issueType: String) : void \\
        }
        
        \umlclass[x=1,y=-6]{TechnicalSupportHandler}{
        }{ 
          + handleRequest(issueType: String) : void \\
        }

        \umlclass[x=9,y=-4]{ManagerSupportHandler}{
        }{ 
          + handleRequest(issueType: String) : void \\
        }   

        \umlclass[x=9,y=-6]{FrontlineSupportHandler}{
        }{ 
          + handleRequest(issueType: String) : void \\
        }                
        
        \umlinherit{BillingSupportHandler}{SupportHandler}
        \umlinherit{TechnicalSupportHandler}{SupportHandler}
        \umlinherit{ManagerSupportHandler}{SupportHandler}
        \umlinherit{FrontlineSupportHandler}{SupportHandler}
        
        \umluniassoc{Client}{SupportHandler}
        \umluniassoc[arg=Nachfolger, angle1=-90, angle2=-149, loopsize=2cm]{SupportHandler}{SupportHandler}
    \end{tikzpicture}
\end{frame}

\begin{frame}[fragile]{Bearbeiter}
	\begin{exampleblock}{SupportHandler.java}
 		\begin{lstlisting}
public abstract class SupportHandler {
    protected SupportHandler nextHandler;

    public void setNextHandler(SupportHandler nextHandler) {
        this.nextHandler = nextHandler;
    }

    public abstract void handleRequest(String issueType);
}
 		\end{lstlisting}
 	\end{exampleblock}
\end{frame}

\begin{frame}[fragile]{Konkreter Bearbeiter 1}
	\begin{exampleblock}{FrontlineSupportHandler.java}
 		\begin{lstlisting}
public class FrontlineSupportHandler extends SupportHandler {
    @Override
    public void handleRequest(String issueType) {
        if (issueType.equals("General")) {
            System.out.println("...");
        } else {
            if (nextHandler != null) {
                System.out.println("...");
                nextHandler.handleRequest(issueType);
            }
        }
    }
}    
 		\end{lstlisting}
 	\end{exampleblock}
\end{frame}

\begin{frame}[fragile]{Konkreter Bearbeiter 2}
	\begin{exampleblock}{TechnicalSupportHandler.java}
 		\begin{lstlisting}
public class TechnicalSupportHandler extends SupportHandler {
    @Override
    public void handleRequest(String issueType) {
        if (issueType.equals("Technical")) {
            System.out.println("...");
        } else {
            if (nextHandler != null) {
                System.out.println("...");
                nextHandler.handleRequest(issueType);
            }
        }
    }
}
        \end{lstlisting}
 	\end{exampleblock}
\end{frame}

\frameWithImageOnly{\BILDERPFAD/chainOfResponsibility2.jpg}
\note{
    Ein weiteres Beispiel für das Chain of Responsibility-Designmuster ist ein Zugriffskontrollsystem (z.B. für Benutzerrechte). Hier müssen verschiedene Zugriffsebenen überprüft werden, bevor einem Benutzer Zugriff auf eine bestimmte Ressource gewährt wird. Jeder Handler überprüft dabei eine spezifische Zugriffsebene und leitet die Anfrage weiter, wenn der aktuelle Zugriff nicht ausreicht.
}

\begin{frame}{Klassendiagramm Zugriffskontrollsystem}
    \begin{tikzpicture}[scale=0.7, transform shape]
        % Define the classes
        \umlclass[x=-3,y=0]{Client}{
        }{ 
        }
        
        \umlclass[x=4,y=0]{AccessHandler}{
        }{ 
          + handleRequest(accessLevel: String) : void \\
        }
        
        \umlclass[x=4,y=-6]{AdminAccessHandler}{
        }{ 
          + handleRequest(accessLevel: String) : void \\
        }
        
        \umlclass[x=8,y=-3]{UserAccessHandler}{
        }{ 
          + handleRequest(accessLevel: String) : void \\
        }   

        \umlclass[x=-1,y=-3]{GuestAccessHandler}{
        }{ 
          + handleRequest(accessLevel: String) : void \\
        }           
        
        \umlinherit{GuestAccessHandler}{AccessHandler}
        \umlinherit{AdminAccessHandler}{AccessHandler}
        \umlinherit{UserAccessHandler}{AccessHandler}
        
        \umluniassoc{Client}{AccessHandler}
        \umluniassoc[arg=Nachfolger, angle1=-90, angle2=-149, loopsize=2cm]{AccessHandler}{AccessHandler}
    \end{tikzpicture}
\end{frame}

\begin{frame}[fragile]{Bearbeiter.java}
	\begin{exampleblock}{AccessHandler}
 		\begin{lstlisting}
public abstract class AccessHandler {
    protected AccessHandler nextHandler;

    public void setNextHandler(AccessHandler nextHandler) {
        this.nextHandler = nextHandler;
    }

    public abstract void handleRequest(String accessLevel);
}
 		\end{lstlisting}
 	\end{exampleblock}
\end{frame}

\begin{frame}[fragile]{Konkreter Bearbeiter 1}
	\begin{exampleblock}{GuestAccessHandler.java}
 		\begin{lstlisting}
public class GuestAccessHandler extends AccessHandler {
    @Override
    public void handleRequest(String accessLevel) {
        if (accessLevel.equals("Guest")) {
            System.out.println("...");
        } else {
            if (nextHandler != null) {
                System.out.println("G...");
                nextHandler.handleRequest(accessLevel);
            }
        }
    }
}
 		\end{lstlisting}
 	\end{exampleblock}
\end{frame}

\begin{frame}[fragile]{Konkreter Bearbeiter 2}
	\begin{exampleblock}{UserAccessHandler.java}
 		\begin{lstlisting}
public class UserAccessHandler extends AccessHandler {
    @Override
    public void handleRequest(String accessLevel) {
        if (accessLevel.equals("User")) {
            System.out.println("U...");
        } else {
            if (nextHandler != null) {
                System.out.println("...");
                nextHandler.handleRequest(accessLevel);
            }
        }
    }
}
        \end{lstlisting}
 	\end{exampleblock}
\end{frame}

\frameWithImageOnly{\BILDERPFAD/chainOfResponsibility3.jpg}
\note{
    Das letzte Beispiel ist ein Logging-System mit verschiedenen Log-Stufen ähnlich wie bei Log4J.
}

\begin{frame}{Klassendiagramm Logsystem}
    \begin{tikzpicture}[scale=0.7, transform shape]
        % Define the classes
        \umlclass[x=-3,y=0]{Client}{
        }{ 
        }
        
        \umlclass[x=4,y=0]{AbstractLogger}{
        }{ 
          + writeMessage(message: String) : void \\
        }
        
        \umlclass[x=4,y=-6]{ErrorLogger}{
        }{ 
          + writeMessage(message: String) : void \\
        }
        
        \umlclass[x=0,y=-3]{FileLogger}{
        }{ 
          + writeMessage(message: String) : void \\
        }   

        \umlclass[x=8,y=-3]{ConsoleLogger}{
        }{ 
          + writeMessage(message: String) : void \\
        }           
        
        \umlinherit{ConsoleLogger}{AbstractLogger}
        \umlinherit{ErrorLogger}{AbstractLogger}
        \umlinherit{FileLogger}{AbstractLogger}
        
        \umluniassoc{Client}{AbstractLogger}
        \umluniassoc[arg=Nachfolger, angle1=-90, angle2=-149, loopsize=2cm]{AbstractLogger}{AbstractLogger}
    \end{tikzpicture}
\end{frame}

\begin{frame}[fragile]{Bearbeiter}
	\begin{exampleblock}{AbstractLogger.java}
 		\begin{lstlisting}
public abstract class AbstractLogger {
	private LogLevel logLevel;
	private AbstractLogger nextLogger;
	
	public AbstractLogger(LogLevel logLevel, 
			AbstractLogger nextLogger) {
		super();
		this.logLevel = logLevel;
		this.nextLogger = nextLogger;
	}
	
	public AbstractLogger(LogLevel logLevel) {
		super();
		this.logLevel = logLevel;
	}
 		\end{lstlisting}
 	\end{exampleblock}
\end{frame}

\begin{frame}[fragile]{Bearbeiter}
	\begin{exampleblock}{AbstractLogger.java}
 		\begin{lstlisting}
	public void logMessage(LogLevel logLevel, String message) {
		if (this.logLevel.ordinal() <= logLevel.ordinal()) {
			writeMessage(message);
		}
		
		if (nextLogger != null) {
			nextLogger.logMessage(logLevel, message);
		}
	}
	
	public abstract void writeMessage(String message);
	
	protected AbstractLogger getNextLogger() {
		return nextLogger;
	}
} 		    
 		\end{lstlisting}
 	\end{exampleblock}
\end{frame}

\begin{frame}[fragile]{Konkreter Bearbeiter}
	\begin{exampleblock}{ConsoleLogger.java}
 		\begin{lstlisting}
public class ConsoleLogger extends AbstractLogger {
	public ConsoleLogger(LogLevel logLevel, 
			AbstractLogger nextLogger) {
		super(logLevel, nextLogger);
	}
	
	public ConsoleLogger(LogLevel logLevel) {
		super(logLevel);
	}
	
	@Override
	public void writeMessage(String message) {
		System.out.println("Standard Console: " + message);
	}
}
        \end{lstlisting}
 	\end{exampleblock}
\end{frame}

\begin{frame}[fragile]{Klient}
	\begin{exampleblock}{LoggerMain.java}
 		\begin{lstlisting}
private static AbstractLogger getChainOfLoggers() {
    AbstractLogger errorLogger = new ErrorLogger(
            LogLevel.ERROR);
    AbstractLogger fileLogger = new FileLogger(
            LogLevel.DEBUG, errorLogger);
    AbstractLogger consoleLogger = new ConsoleLogger(
            LogLevel.INFO, fileLogger);
    
    return consoleLogger;
}
 		\end{lstlisting}
 	\end{exampleblock}
\end{frame}